%% ============================================================================
\section{Introduction}
\label{sec:intro}
\label{sec:voyager}
%% ============================================================================

The dominant paradigm for evaluating the factuality of long-form generation is \emph{decompose-then-verify}: 
a response is decomposed into atomic claims, each claim is verified against retrieved evidence, 
and the resulting verdicts are aggregated into a final score
\citep{factscore,safe,veriscore,factreasoner}. This framework evaluates claims in isolation, and so discards the relations that make a set of
statements an argument. 

For example, Figure~\ref{fig:voyager} shows two answers about Voyager~1 that state the
same five claims with the same stance, so they receive the same verdicts. Response A derives its conclusion
from its evidence and uses both to refute two misconceptions. Response B sets two of its own truths against
each other and derives a falsehood from a truth. No per-claim check separates them, and
\citet{lorefact} show that factuality pipelines miss such errors.

Discourse coherence does not capture the distinction: it concerns how a text is organized, largely
ignores the truth of its claims, and changes when sentences are reordered \citep{grosz1995,barzilay}.
\emph{Logical coherence} is instead a property of an argument: the extent to which the
truth-functional relations it draws among a fixed set of claims are supported by their truth values.
It is invariant to reordering, and a contradiction can be rhetorically well integrated yet logically
incoherent.

% ---------------------------------------------------------------------------
% The Voyager 1 worked pair (Sec. sec:voyager), modelled on the ICLR motivating
% figure and rescored with the committed-world score. Compact (~1/2 page): per
% response, an abridged text, its relation graph and its per-relation
% contributions Delta_r; one score line below. The full responses are in
% App. app:voyager; the reading of the figure is in Secs. sec:voyager and sec:score.
% Relations transcribe data/lcs/example-7-coherence-pair.json; every number is a
% macro from numbers.tex (scripts/arr_coherence_eval.py, exact enumeration).
% ---------------------------------------------------------------------------
\begin{figure*}[t]
\centering
\footnotesize
\newcommand{\mtvx}{11.5mm}
\newcommand{\mtvy}{9.2mm}
\newcommand{\mtvbox}{21mm}
\newcommand{\dscale}{12}

{\scriptsize\textbf{Query}\ \ \emph{Where is Voyager~1 now, and has it left the Solar System?}}

\smallskip
\noindent
%% ===== Response A ========================================================
\begin{minipage}[t]{0.495\textwidth}\centering
  \colorbox{blue!4}{\parbox[t][\mtvbox][t]{\dimexpr\linewidth-2\fboxsep\relax}{%
  \textbf{\textcolor{blue!55!black}{Response A}}\\[1pt]
  \scriptsize
  Voyager~1 crossed the heliopause in August 2012,\,\claimref{1} and its instruments then
  measured a sharp rise in plasma density\,\claimref{2} --- \emph{the measurement that
  established the boundary had been passed.} Voyager~1 is \emph{therefore} in interstellar
  space.\,\claimref{3} \emph{One misconception is that} the plasma wave instrument had failed
  before the crossing;\,\claimref{4} \emph{it cannot have, since that instrument returned the
  density profile.} \emph{Another is that} it has left the Solar System;\,\claimref{5}
  \emph{being in interstellar space is not the same thing.}\strut}}

  \smallskip
  \begin{tikzpicture}[x=\mtvx, y=\mtvy]
    \node[atomtrue]  (a1) at (0,0)        {$a_1$};
    \node[atomtrue]  (a2) at (1.15,0)     {$a_2$};
    \node[atomtrue]  (a3) at (2.3,0)      {$a_3$};
    \node[atomfalse] (a4) at (1.7,-1.5)   {$a_4$};
    \node[atomfalse] (a5) at (2.9,-1.5)   {$a_5$};
    \node[prlab,above=0.3mm of a1] {$\xs{=}1$};
    \node[prlab,above=0.3mm of a2] {$\xs{=}1$};
    \node[prlab,above=0.3mm of a3] {$\xs{=}1$};
    \node[prlab,below=0.3mm of a4] {$\xs{=}0$};
    \node[prlab,below=0.3mm of a5] {$\xs{=}0$};
    \draw[ent,relw=0.90] (a1) -- (a2) node[wlab,above]{$.90$};
    \draw[ent,relw=0.85] (a2) -- (a3) node[wlab,above]{$.85$};
    \draw[con,relw=0.90] (a3) -- (a4) node[wlab,left,pos=0.5]{$.90$};
    \draw[con,relw=0.85] (a3) -- (a5) node[wlab,right,pos=0.5]{$.85$};
    \path (-0.5,-2.25) -- (3.6,0.75);
  \end{tikzpicture}

  \begin{tikzpicture}[x=1mm, y=1mm]
    \def\zx{34}\def\rh{3.4}
    \node[font=\scriptsize\bfseries,anchor=west] at (0,2.6) {$\Delta_r$ (change in $\Lambda$)};
    \foreach \i/\lab/\v in {%
        1/{$a_1\Rightarrow a_2$}/\vyADRawOne,
        2/{$a_2\Rightarrow a_3$}/\vyADRawTwo,
        3/{$a_3\ \con\ a_4$}/\vyADRawThree,
        4/{$a_3\ \con\ a_5$}/\vyADRawFour}
    {
      \pgfmathsetmacro{\y}{-\i*\rh}
      \pgfmathsetmacro{\w}{\v*\dscale}
      \node[anchor=west,font=\scriptsize] at (0,\y) {\lab};
      \fill[blue!35!white,draw=blue!55!black,line width=0.3pt]
           (\zx,\y-1.1) rectangle (\zx+\w,\y+1.1);
      \node[anchor=west,font=\tiny] at (\zx+\w+0.8,\y)
           {$+\pgfmathprintnumber[fixed,fixed zerofill,precision=2]{\v}$};
    }
    \draw[black!60,line width=0.4pt] (\zx,-0.8) -- (\zx,-6*\rh-1.6);
    \path (0,-6*\rh-1.6) -- (60,-6*\rh-1.6);
  \end{tikzpicture}
\end{minipage}\hfill
%% ===== Response B ========================================================
\begin{minipage}[t]{0.495\textwidth}\centering
  \colorbox{red!4}{\parbox[t][\mtvbox][t]{\dimexpr\linewidth-2\fboxsep\relax}{%
  \textbf{\textcolor{red!60!black}{Response B}}\\[1pt]
  \scriptsize
  Voyager~1 crossed the heliopause in August 2012,\,\claimref{1} and its instruments then
  measured a sharp rise in plasma density.\,\claimref{2} \emph{That said, the probe cannot
  really be counted as being in interstellar space.}\,\claimref{3} The plasma wave instrument
  had in fact failed before the crossing,\,\claimref{4} \emph{and that failure tells us the
  probe reached interstellar space;} \emph{it also means} Voyager~1 has left the Solar
  System.\,\claimref{5} \emph{Of course, the crossing rules out any such failure, and the
  density rise rules out having left.}\strut}}

  \smallskip
  \begin{tikzpicture}[x=\mtvx, y=\mtvy]
    \node[atomtrue]  (a1) at (0,0)        {$a_1$};
    \node[atomtrue]  (a2) at (1.15,0)     {$a_2$};
    \node[atomtrue]  (a3) at (2.3,0)      {$a_3$};
    \node[atomfalse] (a4) at (0.55,-1.5)  {$a_4$};
    \node[atomfalse] (a5) at (1.75,-1.5)  {$a_5$};
    \node[prlab,above=0.3mm of a1] {$\xs{=}1$};
    \node[prlab,above=0.3mm of a2] {$\xs{=}1$};
    \node[prlab,above=0.3mm of a3] {$\xs{=}1$};
    \node[prlab,left=0.4mm of a4]  {$\xs{=}0$};
    \node[prlab,right=0.4mm of a5] {$\xs{=}0$};
    \draw[ent,relw=0.90] (a1) -- (a2) node[wlab,above]{$.90$};
    \draw[con,relw=0.90] (a2) -- (a3) node[wlab,above]{$.90$};
    \draw[ent,relw=0.88] (a4) -- (a3)
          node[wlab,above left=-0.4mm and -1.0mm,pos=0.30]{$.88$};
    \draw[con,relw=0.90] (a1) -- (a4) node[wlab,left,pos=0.55]{$.90$};
    \draw[ent,relw=0.85] (a4) -- (a5) node[wlab,below]{$.85$};
    \draw[con,relw=0.88] (a2) to[out=-8,in=25]
          node[wlab,right,pos=0.55,inner sep=1pt]{$.88$} (a5);
    \node[chgnum] (n1) at (3.0,-0.35) {1};
    \draw[chglead] (n1) -- (1.85,-0.06);
    \node[chglab,anchor=west] at (3.12,-0.35) {denies a truth};
    \begin{scope}[on background layer]
      \node[chg,circle,minimum size=8.0mm,inner sep=0pt] at (0.55,-1.5) {};
    \end{scope}
    \node[chgnum] (n2) at (-0.45,-0.85) {2};
    \draw[chglead] (n2) -- (0.25,-1.3);
    \node[chglab,anchor=north,align=center] at (1.15,-2.0)
         {premise from a falsehood};
    \path (-0.5,-2.25) -- (3.6,0.75);
  \end{tikzpicture}

  \begin{tikzpicture}[x=1mm, y=1mm]
    \def\zx{44}\def\rh{3.4}
    \node[font=\scriptsize\bfseries,anchor=west] at (0,2.6) {$\Delta_r$ (change in $\Lambda$)};
    \foreach \i/\lab/\v/\neg in {%
        1/{$a_1\Rightarrow a_2$}/\vyBDRawOne/0,
        2/{$a_2\ \con\ a_3$ \textcolor{orange!60!black}{\textbf{1}}}/\vyBDRawTwo/1,
        3/{$a_4\Rightarrow a_3$ \textcolor{orange!60!black}{\textbf{2}}}/\vyBDRawThree/0,
        4/{$a_4\Rightarrow a_5$}/\vyBDRawFour/0,
        5/{$a_1\ \con\ a_4$}/\vyBDRawFive/0,
        6/{$a_2\ \con\ a_5$}/\vyBDRawSix/0}
    {
      \pgfmathsetmacro{\y}{-\i*\rh}
      \pgfmathsetmacro{\w}{\v*\dscale}
      \node[anchor=west,font=\scriptsize] at (0,\y) {\lab};
      \ifnum\neg=1
        \fill[red!30!white,draw=red!60!black,line width=0.3pt]
             (\zx+\w,\y-1.1) rectangle (\zx,\y+1.1);
        \node[anchor=west,font=\tiny] at (\zx+0.8,\y)
             {$\pgfmathprintnumber[fixed,fixed zerofill,precision=2]{\v}$};
      \else
        \fill[blue!35!white,draw=blue!55!black,line width=0.3pt]
             (\zx,\y-1.1) rectangle (\zx+\w,\y+1.1);
        \node[anchor=west,font=\tiny] at (\zx+\w+0.8,\y)
             {$+\pgfmathprintnumber[fixed,fixed zerofill,precision=2]{\v}$};
      \fi
    }
    \draw[black!60,line width=0.4pt] (\zx,-0.8) -- (\zx,-6*\rh-1.6);
  \end{tikzpicture}
\end{minipage}

\smallskip
\noindent\rule{\textwidth}{0.4pt}\\[1pt]
{\small\textbf{Logical coherence score} (verdict polarity):\quad
 A:\ $\LCS=\textcolor{blue!55!black}{\mathbf{\vyALCS}}$
 ($\Lambda=\vyALam$)
 \qquad$\longrightarrow$\qquad
 B:\ $\LCS=\textcolor{red!60!black}{\mathbf{\vyBLCS}}$
 ($\Lambda=\vyBLam$)}

\caption{\textbf{Two responses, the same five claims, different coherence.} Responses are
  abridged (full text in App.~\ref{app:voyager}); each claim is labelled with its committed value
  (blue true, red false). Graphs: solid blue \ent{}, dashed red \con{}, width growing with the
  weight. Bars: each relation's leave-one-out contribution $\Delta_r$ to the log-ratio $\Lambda$ (\S\ref{sec:voyager},
  \S\ref{sec:score}).}
\label{fig:voyager}
\end{figure*}


We represent an argument as a Markov random field (MRF) over its claims whose factors encode its
relations. A factuality pipeline's verdicts define a \emph{committed world}, and the \emph{Logical
Coherence Score} (LCS) is the weight of evidence \citep{good1950} the relations give that world. 
Relations the committed world supports raise it and relations it violates lower it: refuting a
falsehood increases coherence, deriving a falsehood from a truth decreases it. LCS asks whether the
argument's relations accord with the facts, not whether the claims are true, so it complements
factuality metrics. %We further show that LCS is strictly monotone under every coherence-improving edit. 
In summary, our contributions are:

\begin{itemize}[leftmargin=1.2em,itemsep=1pt,topsep=2pt]
\item \textbf{A score with guarantees, and its linearization} (\S\ref{sec:framework}). LCS is strictly
monotone under every coherence-improving edit, for any priors and relation graph; within product MRFs
read by a single-world log-ratio this forces sign-definite tables, two-valued if the guarantee must
hold whatever the verdicts, and the committed world as the readout (App.~\ref{app:why}). It splits
into per-relation evidence and an interaction term, and its first-order expansion, the linear
evidence, charges violations more than it credits supports (\S\ref{sec:decomp}).
\item \textbf{A controlled evaluation on LoCoBench} (\S\ref{sec:experiments}). LCS rises on
\sThreeGraphGptLamIncPct\% of the declared increases with GPT-6-mined graphs. Once ties count as half, the data separate neither LCS, the linear evidence nor a plain signed sum;
we recommend the linear evidence on theoretical grounds. Judges rating one
response at a time do not see added valid inferences; a pairwise judge given the verdicts trails LCS
slightly on increases and keeps far more invariances.
\item \textbf{An evaluation on natural text} (\S\ref{sec:exp-external}). Read with every claim taken as asserted, scores of mined relations fall when a detected contradiction
is added, track essay quality beyond length about as well as a relation count, and relate to rated
discourse coherence only through which relations are mined.
\end{itemize}
